\Askhsh{13442}{2}
\begin{ylh}{1}
    \item 3.1. Είδη και στοιχεία τριγώνων
    \item 3.2. 1ο Κριτήριο ισότητας τριγώνων
    \item 4.6. Άθροισμα γωνιών τριγώνου.
\end{ylh}
% ===================================================================
%  Αρχείο: 13442.tex (Fragment)
%  Αυτόματη μετατροπή μέσω doc2latex.py
% ===================================================================

ΘΕΜΑ 2

Δίνεται ορθογώνιο τρίγωνο ΑΒΓ με $\widehat{Α} = 90^{ο}$. Στην πλευρά του ΑΒ θεωρούμε σημείο Δ ώστε $Β\Delta = \Delta\Gamma$ και $Α\Delta = Α\Gamma$.

\begin{alist}
\item Να αποδείξετε ότι $Α\widehat{\Delta}\Gamma = 45^{ο}$ . \monades{12}

\item Να υπολογίσετε τη γωνία $\widehat{Β}$. \monades{13}


\begin{center}
\begin{tikzpicture}[scale=1.2]
\tkzDefPoint(0,0){A}
\tkzDefPoint(2.5,0){C}
\tkzDefPoint(0,2.5){D}
% BD = DC. Let D = (0, AC). Then DC^2 = AC^2 + AD^2 = 2 AC^2.
% BD = DC => BD^2 = 2 AC^2. So BD = sqrt(2) * AC.
\tkzDefPoint(0, {2.5 + 2.5*sqrt(2)}){B}
\draw[l3] (A)--(B)--(C)--cycle;
\draw[l3] (D)--(C);
\tkzDrawPoints(A,B,C,D)
\tkzLabelPoint[below left](A){$A$}
\tkzLabelPoint[above](B){$B$}
\tkzLabelPoint[right](C){$\Gamma$}
\tkzLabelPoint[left](D){$\Delta$}
\tkzMarkRightAngle(B,A,C)
\tkzMarkSegments[mark=s|](A,D A,C)
\tkzMarkSegments[mark=s||](B,D D,C)
\end{tikzpicture}
\end{center}
\end{alist}
